% ============================================================
% AL2002 Final Project Report
% Intelligent Reading Comprehension & Quiz Generation System
% ============================================================
\documentclass[12pt,a4paper]{article}

\usepackage[margin=2.5cm]{geometry}
\usepackage{booktabs}
\usepackage{graphicx}
\usepackage{caption}
\usepackage{subcaption}
\usepackage{hyperref}
\usepackage{amsmath}
\usepackage{xcolor}
\usepackage{listings}
\usepackage{float}
\usepackage{array}
\usepackage{multirow}
\usepackage{setspace}
\onehalfspacing

\definecolor{codeblue}{rgb}{0.0,0.3,0.6}
\lstset{
  basicstyle=\ttfamily\small,
  keywordstyle=\color{codeblue}\bfseries,
  commentstyle=\color{gray},
  stringstyle=\color{orange!70!black},
  breaklines=true,
  frame=single,
  numbers=left,
  numberstyle=\tiny\color{gray},
}

\title{
  \textbf{Intelligent Reading Comprehension \&\\Quiz Generation System}\\[0.5em]
  \large AL2002 Final Project Report
}
\author{Zakin Nabeel}
\date{May 2026}

\begin{document}

\maketitle
\tableofcontents
\newpage

% ============================================================
\section{Project Overview}
% ============================================================

This project implements a classical machine-learning pipeline for reading
comprehension and automatic multiple-choice question (MCQ) generation using
the RACE dataset~\cite{lai2017race}.
The system consists of two cooperating models:

\begin{itemize}
  \item \textbf{Model A} — Answer verification: given a passage, a question,
        and four candidate options, predict which option is correct.
  \item \textbf{Model B} — Distractor and hint generation: given a passage
        and the correct answer, generate three plausible wrong options
        (distractors) and ranked passage-level hints.
\end{itemize}

A Streamlit web UI ties both models into an interactive quiz application.
All modelling is done with scikit-learn; no deep learning or GPU is required.

% ============================================================
\section{Dataset}
% ============================================================

\subsection{RACE Dataset}

RACE (Reading Comprehension from Examinations) is collected from Chinese
middle-school and high-school English exams~\cite{lai2017race}.

\begin{table}[H]
\centering
\caption{RACE Dataset Statistics}
\begin{tabular}{lrr}
  \toprule
  \textbf{Property} & \textbf{Value} \\
  \midrule
  Total questions       & $\sim$88,000 \\
  Unique passages       & $\sim$28,000 \\
  Options per question  & 4 (A, B, C, D) \\
  Answer type           & Single correct option \\
  Source                & Chinese school English exams \\
  Avg.\ article length  & 275 words \\
  \bottomrule
\end{tabular}
\end{table}

\subsection{Data Splits}

\begin{table}[H]
\centering
\caption{Dataset Splits Used in Training}
\begin{tabular}{lrrr}
  \toprule
  \textbf{Split} & \textbf{Rows} & \textbf{Unique Articles} & \textbf{Avg.\ Article Words} \\
  \midrule
  Train & 87,866 & 25,135 & 275.0 \\
  Dev   & 87,866 & 25,135 & 275.0 \\
  Test  & 87,866 & 25,135 & 275.0 \\
  \bottomrule
\end{tabular}
\end{table}

\subsection{Class Imbalance}

Each MCQ produces four binary rows (one correct, three incorrect), giving a
3:1 class imbalance ($\text{label}=0 : \text{label}=1 = 15000 : 5000$).
All supervised classifiers use \texttt{class\_weight=`balanced'} to compensate.

% ============================================================
\section{Exploratory Data Analysis}
% ============================================================

\subsection{Answer Distribution}

\begin{figure}[H]
  \centering
  % Replace with actual screenshot from notebook cell eda_answer_dist
  \fbox{\includegraphics[width=0.85\textwidth]{screenshots/eda_answer_dist.png}}
  \caption{Answer label distribution across train, dev, and test splits.
           Labels A--D are approximately balanced ($\sim$25\% each).}
  \label{fig:answer_dist}
\end{figure}

\subsection{Article and Question Length}

\begin{figure}[H]
  \centering
  \fbox{\includegraphics[width=0.85\textwidth]{screenshots/eda_text_len.png}}
  \caption{Distribution of article lengths (left) and question lengths (right)
           in the training split.  The red dashed line marks the mean.}
  \label{fig:text_len}
\end{figure}

\subsection{Option Length: Correct vs Incorrect}

\begin{figure}[H]
  \centering
  \fbox{\includegraphics[width=0.55\textwidth]{screenshots/eda_option_len.png}}
  \caption{Box-plot of option word count for correct (1) vs incorrect (0) options.
           Correct options tend to be slightly longer, providing a weak signal.}
  \label{fig:option_len}
\end{figure}

% ============================================================
\section{Statistical Analysis}
% ============================================================

\subsection{Cosine Similarity as a Discriminating Feature}

We computed TF-IDF cosine similarity between each answer option and the
passage for 500 training samples.

\begin{figure}[H]
  \centering
  \fbox{\includegraphics[width=0.75\textwidth]{screenshots/stat_cosine.png}}
  \caption{Distribution of cosine similarity (option vs.\ article).
           Correct options (green) have higher similarity than incorrect options (red).}
  \label{fig:cosine_dist}
\end{figure}

\subsection{Statistical Significance (t-test)}

\begin{table}[H]
\centering
\caption{Independent two-sample t-test: correct vs incorrect option cosine similarity}
\begin{tabular}{lr}
  \toprule
  \textbf{Statistic} & \textbf{Value} \\
  \midrule
  t-statistic & 2.4300 \\
  p-value     & $1.52 \times 10^{-2}$ \\
  Significant ($p < 0.05$)? & \textbf{Yes} \\
  \bottomrule
\end{tabular}
\end{table}

\textbf{Interpretation:} TF-IDF cosine similarity to the passage is a
statistically significant discriminating feature for answer verification.

% ============================================================
\section{Feature Engineering}
% ============================================================

Each sample (article, question, one option) is represented as a concatenation
of:

\begin{enumerate}
  \item \textbf{TF-IDF vector} (15,000 features, unigrams + bigrams) of the
        combined text \texttt{[article $\times$2, question, option]}.
        The article is repeated to up-weight passage vocabulary.
  \item \textbf{Three cosine similarity scores}:
        \begin{itemize}
          \item $\text{sim}(\text{question}, \text{option})$
          \item $\text{sim}(\text{article}, \text{option})$ — the key signal
          \item $\text{sim}(\text{article}, \text{question})$
        \end{itemize}
\end{enumerate}

Final feature matrix shape: $(N, 15003)$ — 15,000 TF-IDF + 3 cosine.

% ============================================================
\section{Model A: Answer Verification}
% ============================================================

\subsection{Model Selection Rationale}

\begin{table}[H]
\centering
\caption{Model A — Model Selection Rationale}
\begin{tabular}{lp{9cm}}
  \toprule
  \textbf{Model} & \textbf{Why chosen} \\
  \midrule
  Logistic Regression  & Fast, probabilistic (calibrated confidence scores), linear on TF-IDF \\
  Linear SVM           & Strong generalisation on high-dim sparse text, handles imbalance well \\
  Label Propagation    & Semi-supervised: exploits unlabelled examples via k-NN label spreading \\
  KMeans ($k=2$)       & Unsupervised baseline: clusters options without labels \\
  GMM (diag)           & Probabilistic unsupervised clustering with soft cluster assignments \\
  \textbf{MLP Shallow} & Neural network baseline: Dense(256, ReLU) to test non-linear gain \\
  \textbf{MLP Deep}    & Deeper neural net: Dense(512)→Dense(256)→Dense(128) \\
  \bottomrule
\end{tabular}
\end{table}

\subsection{Training Details}

\begin{itemize}
  \item \textbf{Training samples:} 5,000 MCQs $\to$ 20,000 binary rows
  \item \textbf{Logistic Regression:} GridSearchCV over $C \in \{0.1, 1, 5\}$,
        3-fold CV; best $C = 0.1$, CV F1 = 0.5167
  \item \textbf{SVM:} \texttt{LinearSVC}, \texttt{class\_weight=`balanced'}
  \item \textbf{Label Propagation:} 3,000-sample subset, 30\% labels masked,
        \texttt{kernel=`knn'}, \texttt{n\_neighbors=7}
  \item \textbf{KMeans/GMM:} 3,000-sample dense subset, \texttt{n\_components=2}
  \item \textbf{MLP:} Adam optimiser, early stopping, 10\% validation fraction
\end{itemize}

\subsection{Evaluation Results}

\begin{table}[H]
\centering
\caption{Model A — Performance on Dev Set (4,000 samples)}
\begin{tabular}{lrrl}
  \toprule
  \textbf{Model} & \textbf{Accuracy} & \textbf{Macro F1} & \textbf{Type} \\
  \midrule
  Logistic Regression     & 0.5733 & 0.5222 & Classical ML \\
  Linear SVM              & 0.6040 & 0.5594 & Classical ML \\
  Label Propagation       & 0.7771 & 0.5904 & Semi-supervised \\
  MLP Shallow (256)       & 0.6125 & 0.5681 & Neural Network \\
  MLP Deep (512-256-128)  & 0.6083 & 0.5643 & Neural Network \\
  KMeans ($k=2$)          & N/A    & 0.0037$^\dagger$ & Unsupervised \\
  GMM (diag, $k=2$)       & N/A    & 0.0037$^\dagger$ & Unsupervised \\
  \midrule
  Random baseline         & 0.2500 & 0.2500 & — \\
  \bottomrule
  \multicolumn{4}{l}{$\dagger$ Silhouette score (not F1).}
\end{tabular}
\end{table}

\begin{figure}[H]
  \centering
  \fbox{\includegraphics[width=0.85\textwidth]{screenshots/bar_compare.png}}
  \caption{Accuracy and Macro F1 for classical ML models.}
  \label{fig:bar_compare}
\end{figure}

\begin{figure}[H]
  \centering
  \fbox{\includegraphics[width=0.85\textwidth]{screenshots/nn_vs_classical.png}}
  \caption{Classical ML vs Neural Network comparison.  MLP Shallow marginally
           outperforms Linear SVM, but the gap is $<1\%$.}
  \label{fig:nn_compare}
\end{figure}

\subsection{Confusion Matrices}

\begin{figure}[H]
  \centering
  \fbox{\includegraphics[width=0.85\textwidth]{screenshots/confusion_matrices.png}}
  \caption{Confusion matrices for Logistic Regression (left) and Linear SVM (right).
           Both models recall $\sim$50--57\% of correct options (class 1).}
  \label{fig:confusion}
\end{figure}

\subsection{ROC Curve}

\begin{figure}[H]
  \centering
  \fbox{\includegraphics[width=0.55\textwidth]{screenshots/roc_curve.png}}
  \caption{ROC curve for Logistic Regression (AUC = 0.5628).
           The curve above the diagonal confirms the model learns a signal
           beyond random guessing.}
  \label{fig:roc}
\end{figure}

\subsection{Neural Network vs Classical ML — Analysis}

\begin{table}[H]
\centering
\caption{Neural Network vs Classical ML — Key Findings}
\begin{tabular}{p{4.5cm}p{9cm}}
  \toprule
  \textbf{Finding} & \textbf{Explanation} \\
  \midrule
  MLP Shallow $\approx$ SVM & Sparse TF-IDF features have limited non-linear structure;
    a single hidden layer adds $<1\%$ accuracy. \\
  MLP Deep $<$ MLP Shallow & Adding depth hurts: high-dimensional sparse input
    causes overfitting without regularisation. \\
  Label Propagation wins & Semi-supervised learning exploits the unlabelled pool,
    giving a $+17\%$ accuracy advantage over Linear SVM. \\
  All beat random (25\%) & TF-IDF cosine similarity to the passage is a real signal. \\
  No BERT needed here & Transformers would achieve $\sim$70--80\% but require GPU
    and are outside the classical-ML scope of this project. \\
  \bottomrule
\end{tabular}
\end{table}

\subsection{Ensemble Inference}

At inference time, each option receives an ensemble score:
\[
  \text{score}(\text{opt}) = 0.6 \cdot P_{\text{LR}}(\text{correct}) +
  0.4 \cdot \sigma\!\left(f_{\text{SVM}}(\text{opt})\right)
\]
where $\sigma$ is the sigmoid function applied to the SVM decision function.
The option with the highest score is predicted as the answer.

Demo results on 5 random dev samples:
\begin{table}[H]
\centering
\caption{Ensemble Inference Demo (5 samples, seed 99)}
\begin{tabular}{lcccr}
  \toprule
  \textbf{Question (truncated)} & \textbf{Pred.} & \textbf{Gold} & \textbf{Correct?} & \textbf{Latency} \\
  \midrule
  What is important for an English learner? & D & D & \checkmark & 34.5 ms \\
  The best title for this passage \ldots    & B & B & \checkmark & 21.1 ms \\
  If I am driving at 100 km/hr \ldots       & A & D & $\times$ & 18.2 ms \\
  How long did Gauguin live \ldots          & C & A & $\times$ & 18.8 ms \\
  The author thought the Chinese \ldots     & A & B & $\times$ & 19.6 ms \\
  \midrule
  \multicolumn{3}{l}{Demo Accuracy} & \multicolumn{2}{r}{2/5 (40\%)} \\
  \multicolumn{3}{l}{Avg Latency}   & \multicolumn{2}{r}{22.4 ms} \\
  \bottomrule
\end{tabular}
\end{table}

% ============================================================
\section{Model B: Distractor and Hint Generation}
% ============================================================

\subsection{Distractor Generation}

\begin{enumerate}
  \item Extract content-word candidates from the article (excluding stopwords,
        $\text{len} > 3$).
  \item Remove candidates that overlap with the correct answer.
  \item Score each candidate by TF-IDF cosine similarity to the correct answer.
  \item Select the $k$ candidates whose similarity is closest to 0.25
        (plausible but wrong).
\end{enumerate}

\subsection{Hint Generation}

\begin{enumerate}
  \item Split the article into sentences.
  \item Rank sentences by TF-IDF cosine similarity to the question.
  \item Return the top-$k$ sentences in ascending similarity order
        (general $\to$ specific).
\end{enumerate}

\subsection{Evaluation}

\begin{table}[H]
\centering
\caption{Model B — Distractor Evaluation (50 dev samples, 150 distractors)}
\begin{tabular}{lr}
  \toprule
  \textbf{Metric} & \textbf{Value} \\
  \midrule
  Total distractors & 150 \\
  Valid distractors & 150 \\
  Accuracy          & 1.0000 \\
  Precision         & 1.0000 \\
  Recall            & 1.0000 \\
  F1 Score          & 1.0000 \\
  \bottomrule
\end{tabular}
\end{table}

\textbf{Interpretation:} All generated distractors are surface-valid (differ
from the correct answer and contain no placeholder text).  Quality as
perceived by a human reader is harder to quantify without manual annotation.

\subsection{Demo Output}

\begin{figure}[H]
  \centering
  \fbox{\includegraphics[width=0.85\textwidth]{screenshots/model_b_demo.png}}
  \caption{Model B demo: generated distractors and hints for a sample question
           about the Dahab explosions.}
  \label{fig:model_b_demo}
\end{figure}

% ============================================================
\section{Streamlit UI}
% ============================================================

\begin{figure}[H]
  \centering
  \fbox{\includegraphics[width=0.85\textwidth]{screenshots/ui_tab1.png}}
  \caption{Tab 1 — Article Input: load a RACE sample or paste a custom article.}
\end{figure}

\begin{figure}[H]
  \centering
  \fbox{\includegraphics[width=0.85\textwidth]{screenshots/ui_tab2.png}}
  \caption{Tab 2 — Quiz: answer MCQ variants with Model B distractors.}
\end{figure}

\begin{figure}[H]
  \centering
  \fbox{\includegraphics[width=0.85\textwidth]{screenshots/ui_tab3.png}}
  \caption{Tab 3 — Hints: progressively reveal passage-extracted hints.}
\end{figure}

\begin{figure}[H]
  \centering
  \fbox{\includegraphics[width=0.85\textwidth]{screenshots/ui_tab4.png}}
  \caption{Tab 4 — Developer Dashboard: Model A/B metrics, latency chart,
           session history, and CSV export.}
\end{figure}

% ============================================================
\section{System Architecture Summary}
% ============================================================

\begin{table}[H]
\centering
\caption{Full Pipeline — Component Summary}
\begin{tabular}{llp{6.5cm}}
  \toprule
  \textbf{Component} & \textbf{Method} & \textbf{Role} \\
  \midrule
  Preprocessing     & Binary expansion (4 rows/MCQ)                  & Convert MCQ task to binary classification \\
  Features          & TF-IDF (15k) + 3 cosine scores                 & Lexical + semantic similarity \\
  Model A (sup.)    & LR + LinearSVC ensemble (60/40 soft vote)       & Predict correct option \\
  Model A (semi.)   & Label Propagation (knn, 30\% unlabelled)        & Leverage unlabelled data \\
  Model A (unsup.)  & KMeans $k=2$, GMM (diag)                        & Cluster options without labels \\
  NN comparison     & MLP Shallow (256) + MLP Deep (512-256-128)      & Validate classical ML choice \\
  Model B dist.     & TF-IDF cosine $\approx$ 0.25 target             & Generate plausible wrong options \\
  Model B hints     & Sentence ranking by cosine to question          & Extractive passage hints \\
  UI                & Streamlit (4 tabs)                              & Interactive quiz application \\
  \bottomrule
\end{tabular}
\end{table}

% ============================================================
\section{Conclusion}
% ============================================================

\begin{itemize}
  \item TF-IDF cosine similarity to the passage is a statistically significant
        ($p = 0.015$) and practically useful feature for answer verification.
  \item \textbf{Linear SVM} (60.4\% accuracy) is the best classical supervised
        model on this feature space.
  \item \textbf{MLP neural networks} offer marginal improvement ($+0.8\%$) but
        add training complexity; for sparse TF-IDF features, linear classifiers
        are the right choice.
  \item \textbf{Label Propagation} (77.7\% accuracy) demonstrates strong
        semi-supervised performance, outperforming both fully supervised and
        neural network baselines by leveraging unlabelled examples.
  \item \textbf{Model B} generates 100\% surface-valid distractors and
        coherent passage-ranked hints, completing the end-to-end MCQ pipeline.
\end{itemize}

\begin{thebibliography}{9}
  \bibitem{lai2017race}
    Guokun Lai et al.
    \textit{RACE: Large-scale ReAding Comprehension Dataset From Examinations}.
    EMNLP 2017.
    \url{https://aclanthology.org/D17-1082/}
\end{thebibliography}

\end{document}
